\documentclass[12pt]{report}
\usepackage[utf8]{inputenc}
\usepackage[a4paper,margin=2.5cm]{geometry}
\usepackage{graphicx}
\usepackage{natbib}
\usepackage{setspace}
\onehalfspacing

\begin{document}
\tableofcontents

\chapter{Literature Review}
\label{ch:review}

\section{Background}
Early work on soil erosion focused on field plots \citep{smith1990}. Later,
\citet{jones2005} introduced remote sensing methods.

\subsection{Field studies}
Field studies remain the reference method. They are expensive, slow and
limited to small areas, but they measure erosion directly.

\subsection{Remote sensing}
Satellite imagery covers large areas at low cost.
\begin{figure}[htbp]
  \centering
  \includegraphics[width=\linewidth]{ndvi-map.pdf}
  \caption{Vegetation index over the study area.}
  \label{fig:ndvi}
\end{figure}

\section{Research gap}
No study combines both methods over a full watershed, as Chapter~\ref{ch:review}
shows.

\bibliographystyle{plainnat}
\bibliography{references}
\end{document}
